% !TeX root = ../教材题库_全部条目_转写工作稿.tex
% Source: XJ-B1-C06-TH-P248-Y0430
% 第6章 / 6.4 用样本估计总体 / 教材 p.248 / PDF 第255页
% OCR未校对：公式与文字均可能有误；未完成数学TeX转写；禁止直接发给学生
% 原题及相邻原解截图（核对用）：../题目原图/XJ-B1-C06-TH-P248-Y0430-01.webp
% 本文件只在转写完成后作为题干使用；OCR 条目尚不是可靠题干。
\begin{bankitem}{p.248 思考（含紧邻说明）}{第6章 / 6.4 用样本估计总体 / 教材 p.248 / PDF 第255页}{XJ-B1-C06-TH-P248-Y0430}{OCR未校对：公式与文字均可能有误；未完成数学TeX转写；禁止直接发给学生}
\begin{ocrdraft}
[教材p.248 / PDF第255页]
\par
？思考如何根据频率分布直方图，估计出样本的平均数、中位数和众数？
\par
在频率分布直方图中，由于无法知道每个组内的数据是如何分布的，因此我们
\par
通常假设它们在组内均匀分布．
\par
根据平均数的意义可知，样本平均数可以表示为数据与它的频率的乘积之和，
\par
从而在频率分布直方图中，样本平均数可以用每个小矩形底边中点的横坐标与小矩
\par
形的面积的乘积之和近似代替．这样我们就可近似估计样本的平均数．
\par
根据中位数的意义可知，在样本中，有50％的个体小于或等于中位数，也有
\par
50％的个体大于或等于中位数．因此，在频率分布直方图中，中位数左边和右边的
\par
直方图的面积应该相等，由此可估计样本中位数的值．
\par
根据众数的意义可知，可将频率分布直方图中最高矩形所在区间的中点所对应
\par
的数据作为样本众数的估计值．
\par
下面进行举例说明．
\end{ocrdraft}
\end{bankitem}
